\documentclass[10pt]{article}

\usepackage{parskip}
\usepackage[margin=1in]{geometry} 
\usepackage{amsmath,amsthm,amssymb, graphicx, multicol, array}
\usepackage{makecell}
\usepackage{enumitem}
\usepackage{amssymb}
\usepackage{float}
\usepackage{href-ul}
 
\newcommand{\N}{\mathbb{N}}
\newcommand{\Z}{\mathbb{Z}}
 
\newenvironment{problem}[2][Problem]{\begin{trivlist}
\item[\hskip \labelsep {\bfseries #1}\hskip \labelsep {\bfseries #2.}]}{\end{trivlist}}

\begin{document}
 
\title{CSS Network Design}
\author{Jakob Sverre Alexandersen\\
GRA4154 Supply Chain Optimization with Mathematical Programming\\
Lecture 5 (guest)}
\maketitle

\tableofcontents
\newpage

\section{Intro}
GreenCarbon (GC) is planning to build carbon capture and storage (CCS) facilities to capture
and store CO$_2$ emissions from industrial sites. The plan is to build CO$_2$ capture facilities near
industrial emission sources in Antwerp, Milan, Frankfurt, and Bilbao. GC has identified three
potential locations for CO$_2$ storage facilities near Rotterdam, Hamburg, Marseille, and
Copenhagen. Captured CO$_2$ emissions will be transported and stored underground in these
storage facilities. The transport cost is 0.04 EUR per ton CO$_2$ per kilometer.

Table below demonstrates the distances (in km) between industrial emission sources and the
potential CCS facilities
\begin{center}
    \begin{tabular}{|c|c|c|c|c|}
        \hline 
        \textbf{Source} & \textbf{Rotterdam} & \textbf{Hamburg} & \textbf{Marseille} & \textbf{Copenhagen}\\\hline 
        \textbf{Antwerp} & 110 & 460 & 1050 & 900\\\hline
        \textbf{Milan} & 990 & 1100 & 480 & 1300\\\hline
        \textbf{Frankfurt} & 440 & 500 & 750 & 650\\\hline
        \textbf{Bilbao} & 1200 & 1400 & 780 & 1750\\\hline
    \end{tabular}
\end{center}

The following table lists the fixed annual operating costs, maximum annual storage capacities,
and variable storage costs at each potential facility
\begin{center}
    \begin{tabular}{|c|c|c|c|}
        \hline 
        \textbf{Facility} & \textbf{Fixed annual cost (EUR)} & \textbf{Storage Capacity} & \textbf{Variable storage cost}\\\hline 
        \textbf{Rotterdam} & 2,000,000 & 250,000 & 12\\\hline
        \textbf{Hamburg} & 1,800,000 & 200,000 & 11\\\hline
        \textbf{Marseille} & 1,900,000 & 230,000 & 13\\\hline
        \textbf{Copenhagen} & 1,400,000 & 150,000 & 10\\\hline
    \end{tabular}
\end{center}

Table below shows annual CO$_2$ emissions available t for capture from each industrial source: 
\begin{center}
    \begin{tabular}{|c|c|}
        \hline 
        \textbf{Source} & \textbf{Annual Emissions}\\\hline 
        \textbf{Antwerp} & 150,000\\\hline 
        \textbf{Milan} & 130,000\\\hline 
        \textbf{Frankfurt} & 140,000\\\hline 
        \textbf{Bilbao} & 120,000\\\hline 
    \end{tabular}
\end{center}

\section{Model}

\textbf{Sets}
\begin{itemize}
    \item $\mathcal{I} = \{\text{Antwerp, Milan, Frankfurt, Bilbao}\}$: emission sources
    \item $\mathcal{J} = \{\text{Rotterdam, Hamburg, Marseille, Copenhagen}\}$: storage facilities
\end{itemize}

\textbf{Parameters}
\begin{itemize}
    \item $d_{ij}$: distance (km) from source $i $ to facility $j $
    \item $t = 0.04 $: transport cost (EUR/ton/km)
    \item $f_j $: fixed annual operating cost at facility $j $
    \item $s_j $: maximum storage capacity at facility $j $
    \item $v_j $: variable storage cost at facility $j $
    \item $e_i $: annual CO$_2$ emissions available for capture at source $i $
\end{itemize}
\newpage
\textbf{Decision variables}
\begin{align}
    x_{ij} &\ge 0 \quad\forall\, i\in \mathcal{I},\, j\in \mathcal{J}\\
    y_j&\in \{0, 1\}\quad\forall\, j\in J 
\end{align}
\begin{enumerate}
    \item Tons of CO$_2$ transported from source $i $ to facility $j $
    \item 1 if facility $j $ is opened, 0 otherwise 
\end{enumerate}

\textbf{Objective function}
\begin{gather*}
    \min Z = \sum_{j\in \mathcal{J}}f_jy_j + \sum_{i\in \mathcal{I}}\sum_{j\in \mathcal{J}}\left(td_{ij} + v_j\right)x_{ij}\\
    \text{s.t.}\\
    \sum_{j\in \mathcal{J}}x_{ij} = e_i\quad\forall\, i\in \mathcal{I}\\
    \sum_{i\in \mathcal{I}}x_{ij}\le s_jy_j\quad\forall\, j \in \mathcal{J}\\
    x_{ij}\ge 0\quad\forall\, i\in \mathcal{I},\, j\in \mathcal{J}\\
    y_j\in \{0, 1\}\quad\forall\, j\in \mathcal{J}
\end{gather*}



\end{document}